\metatitle{Research with AI tools}
\metadescription{How to use AI tools on the open problems of the symmetric
functions catalog: where to find problems, useful tools, and how to share a
solution and get credit for it.}
\metakeywords{open problems, AI-assisted research, conjectures, Lean,
real-rootedness}

\section[aiResearch]{Research with AI tools}

Do you have AI tokens to spare and want to spend them on mathematics?  This
catalog contains many open problems, and several of them have recently been
solved with substantial help from AI tools; see
\hyperref[aiResearchSolved]{below}.  This page is meant for anyone who wants
to try, in particular students and people outside academia who find
something and want to let me know.

If you want a problem to work on, send me an email (see the
\hyperref[about]{about page}) with a few words about your background and
setup, and I can suggest something that fits.  Otherwise, pick one of the
problems below.

\subsection[aiResearchProblems]{Where to find problems}

Good places to start are problems about one explicit family of polynomials,
where small cases can be computed and the statement can be formalized:
\begin{itemize}
\item \hyperref[realrootedOpenProblems]{Open problems on real-rootedness}:
conjectures with definitions, first polynomials and known partial results,
and a list of counterexamples.
\item \hyperref[personalMinorProblems]{Minor research problems}: smaller
questions related to my own research.
\item \hyperref[gammaPositivityOpenProblems]{Open problems on
gamma-positivity}.
\end{itemize}
Harder, well-known conjectures are stated on the pages on
\hyperref[chromaticEExpansionConjectures]{elementary positivity of chromatic
symmetric functions}, \hyperref[ehrhartConjecturalPositivity]{Ehrhart
positivity}, \hyperref[plethysmProblems]{plethysm},
\hyperref[hallLittlewoodConjectures]{Hall--Littlewood polynomials} and
\hyperref[jackCharacterConjecture]{Jack polynomials}.  Many other pages state
conjectures in passing; the
\href{https://www.symmetricfunctions.com/search.htm}{search} finds them.

\subsection[aiResearchTools]{Tools}

These repositories are public and may be useful.
\begin{itemize}
\item \href{https://github.com/PerAlexandersson/RealRooted}{\icon{github}
RealRooted}: an experimental Lean~4 library on real-rooted polynomials and
interlacing, with many definitions and criteria that a proof can build on.
Pull requests with formalized solutions are welcome.
\item \href{https://github.com/PerAlexandersson/polytool}{\icon{github}
polytool}: exact univariate polynomial tools for real-rootedness, interlacing,
recurrences and sequences, with an MCP server.
\item \href{https://github.com/PerAlexandersson/ehrcalc}{\icon{github}
ehrcalc}: exact Ehrhart computations and related counting tools, with an
\href{https://github.com/PerAlexandersson/ehrcalc/tree/main/mcp}{MCP server}.
\item \href{https://github.com/PerAlexandersson/mcp-oeis}{\icon{github}
mcp-oeis}: an MCP server for searching the OEIS.
\item \href{https://github.com/PerAlexandersson/Mathematica-packages}{\icon{github}
Mathematica packages} for symmetric functions and algebraic combinatorics.
\item \href{https://www.arxiv.symmetricfunctions.com}{arXiv++}: recent
arXiv papers in combinatorics, tagged with topics of this catalog.
\end{itemize}

The catalog itself is written in LaTeX, and its source is on
\href{https://github.com/PerAlexandersson/symmetricfunctions.com}{\icon{github}
GitHub}.

\subsection[aiResearchSharing]{Sharing a solution}

AI tools make mistakes, often in a confident tone, so check every step
yourself before announcing a result.  A formal proof in Lean is the most
convincing format, but make sure that the formal statement really is the
statement of the problem.

\begin{itemize}
\item Put the Lean code in a public GitHub repository, with a short README
saying what is proved and how the statement matches the problem.  A short
informal write-up of the argument is welcome too.  Please use the Apache~2.0
license, as Mathlib and RealRooted do, so that the code can be reused there.
\item Tell me about it, by email or as an
\href{https://github.com/PerAlexandersson/symmetricfunctions.com/issues}{issue
on GitHub}, with the link, the problem it solves, and how you want to be
credited.
\item The result is then credited to you on the corresponding page, with a
link to your repository, and the problem is moved from the open problems to
the relevant topic page.
\item If the proof fits the
\href{https://github.com/PerAlexandersson/RealRooted}{RealRooted} library,
you can also open a pull request there (see its
\href{https://github.com/PerAlexandersson/RealRooted\#contributing}{contributing
notes}); otherwise I may port it myself, again with credit.
\end{itemize}
A public repository also dates the result, which settles priority if
several people find a solution at the same time.  Using AI tools is fine and
there is no need to apologize for it, but please say how they were used.

If you do write a manuscript, I am happy to take a brief look at it.

\subsection[aiResearchPending]{Solutions that are not yet public}

Sometimes I know of a solution before it is public, for instance when a
write-up is being prepared.  Such results are stated on the catalog with a
reference to the unpublished notes, so that nobody spends effort on a problem
that is already settled.

Other people may already be working on the same problem. If you decide to
work on an open problem, consider opening an issue in the
\href{https://github.com/PerAlexandersson/symmetricfunctions.com/issues}{\icon{github}
catalog's GitHub repository} so others know what you are working on.

\subsection[aiResearchSolved]{Recently settled}

Some problems from this catalog that were settled in 2026, several of them
with substantial help from AI tools:
\begin{itemize}
\item \hyperref[conjRibbonMatroidZ]{Narayana $Z$-polynomials of ribbon
lattice path matroids}, by \name{Steve Raj} and, independently, by
ProgramComputer.
\item \hyperref[catalanWordPeaksRealRooted]{Real-rootedness of strict peaks
in Catalan words}, by ProgramComputer.
\item \hyperref[conjLongCycleDescents]{Real-rootedness of descent
polynomials of permutations with long cycles}, by \name{Bowen Zuo}.
\item \hyperref[fussCatalanNNERealRooted]{Real-rootedness of Fuss--Catalan
path polynomials}, by \name{Ryan Halper}.
\item \hyperref[conjTransportationHStar]{A counterexample to real-rootedness
of $h^*$-polynomials of transportation polytopes}, by \name{Blake Stone}.
\end{itemize}
